\documentclass[UTF8,12pt]{ctexart}
\usepackage[paperwidth=240mm,paperheight=200mm,margin=14mm,top=8mm]{geometry}
\usepackage{amsmath,amssymb,mathtools}
\usepackage{xcolor}
\pagestyle{empty}
\setlength{\parindent}{0pt}
\setlength{\parskip}{0pt}
\newcommand{\qn}[1]{\textbf{\large #1}\ }
\xeCJKDeclareCharClass{CJK}{"2460 -> "24FF}
\begin{document}
\qn{2024.1}
函数 $f(x)=|x|^{\frac{1}{(1-x)(x-2)}}$ 的第一类间断点的个数是（　　）
\\
\noindent (A) 3. \qquad (B) 2. \qquad (C) 1. \qquad (D) 0.

\vfill
\newpage
\qn{2024.2}
设函数 $y=f(x)$ 由参数方程 $\begin{cases} x=1+t^3, \\ y=\mathrm{e}^t \end{cases}$ 确定，则 $\lim\limits_{x\to+\infty}x\left[f\left(2+\dfrac{2}{x}\right)-f(2)\right]=$（　　）
\\
\noindent (A) $2\mathrm{e}$. \qquad (B) $\dfrac{4\mathrm{e}}{3}$. \qquad (C) $\dfrac{2\mathrm{e}}{3}$. \qquad (D) $\dfrac{\mathrm{e}}{3}$.

\vfill
\newpage
\qn{2024.3}
设函数 $f(x)=\int_0^{\sin x}\sin t^3\,\mathrm{d}t$，$g(x)=\int_0^x f(t)\,\mathrm{d}t$，则（　　）
\\
\noindent (A) $f(x)$ 是奇函数，$g(x)$ 是奇函数.
\\
\noindent (B) $f(x)$ 是奇函数，$g(x)$ 是偶函数.
\\
\noindent (C) $f(x)$ 是偶函数，$g(x)$ 是偶函数.
\\
\noindent (D) $f(x)$ 是偶函数，$g(x)$ 是奇函数.

\vfill
\newpage
\qn{2024.4}
已知数列 $\{a_n\}$（$a_n\neq0$），若 $\{a_n\}$ 发散，则（　　）
\\
\noindent (A) $\left\{a_n+\dfrac{1}{a_n}\right\}$ 发散.
\\
\noindent (B) $\left\{a_n-\dfrac{1}{a_n}\right\}$ 发散.
\\
\noindent (C) $\left\{\mathrm{e}^{a_n}+\dfrac{1}{\mathrm{e}^{a_n}}\right\}$ 发散.
\\
\noindent (D) $\left\{\mathrm{e}^{a_n}-\dfrac{1}{\mathrm{e}^{a_n}}\right\}$ 发散.

\vfill
\newpage
\qn{2024.5}
已知函数 $f(x,y)=\begin{cases}(x^2+y^2)\cdot\sin\dfrac{1}{xy},&xy\neq0,\\[4pt]0,&xy=0,\end{cases}$ 则在点 $(0,0)$ 处（　　）
\\
\noindent (A) $\dfrac{\partial f(x,y)}{\partial x}$ 连续，$f(x,y)$ 可微.
\\
\noindent (B) $\dfrac{\partial f(x,y)}{\partial x}$ 连续，$f(x,y)$ 不可微.
\\
\noindent (C) $\dfrac{\partial f(x,y)}{\partial x}$ 不连续，$f(x,y)$ 可微.
\\
\noindent (D) $\dfrac{\partial f(x,y)}{\partial x}$ 不连续，$f(x,y)$ 不可微.

\vfill
\newpage
\qn{2024.6}
设 $f(x,y)$ 是连续函数，则 $\int_0^{\frac{\pi}{2}}\mathrm{d}x\int_{\sin x}^1 f(x,y)\,\mathrm{d}y=$（　　）
\\
\noindent (A) $\int_0^1\mathrm{d}y\int_{\frac{\pi}{6}}^{\arcsin y}f(x,y)\,\mathrm{d}x$ \qquad (B) $\int_0^1\mathrm{d}y\int_{\arcsin y}^{\frac{\pi}{2}}f(x,y)\,\mathrm{d}x$
\\
\noindent (C) $\int_0^1\mathrm{d}y\int_0^{\arcsin y}f(x,y)\,\mathrm{d}x$ \qquad (D) $\int_0^1\mathrm{d}y\int_0^{\frac{\pi}{2}}f(x,y)\,\mathrm{d}x$

\vfill
\newpage
\qn{2024.7}
设非负函数 $f(x)$ 在 $[0,+\infty)$ 上连续，给出以下三个命题：
① 若 $\int_0^{+\infty}f^2(x)\,\mathrm{d}x$ 收敛，则 $\int_0^{+\infty}f(x)\,\mathrm{d}x$ 收敛.
② 若存在 $p>1$，使 $\lim\limits_{x\to+\infty}x^p f(x)$ 存在，则 $\int_0^{+\infty}f(x)\,\mathrm{d}x$ 收敛.
③ 若 $\int_0^{+\infty}f(x)\,\mathrm{d}x$ 收敛，则存在 $p>1$，使 $\lim\limits_{x\to+\infty}x^p f(x)$ 存在.
其中真命题的个数为（　　）
\\
\noindent (A) 0. \qquad (B) 1. \qquad (C) 2. \qquad (D) 3.

\vfill
\newpage
\qn{2024.8}
设 $\boldsymbol{A}$ 为3阶矩阵，$\boldsymbol{P}=\begin{pmatrix}1&0&0\\0&1&0\\1&0&1\end{pmatrix}$，若 $\boldsymbol{P}^{\mathrm{T}}\boldsymbol{A}\boldsymbol{P}^2=\begin{pmatrix}a+2c&0&c\\0&b&0\\2c&0&c\end{pmatrix}$，则 $\boldsymbol{A}=$（　　）
\\
\noindent (A) $\begin{pmatrix}c&0&0\\0&a&0\\0&0&b\end{pmatrix}$. \qquad (B) $\begin{pmatrix}b&0&0\\0&c&0\\0&0&a\end{pmatrix}$.
\\
\noindent (C) $\begin{pmatrix}a&0&0\\0&b&0\\0&0&c\end{pmatrix}$. \qquad (D) $\begin{pmatrix}c&0&0\\0&b&0\\0&0&a\end{pmatrix}$.

\vfill
\newpage
\qn{2024.9}
设 $\boldsymbol{A}$ 为4阶矩阵，$\boldsymbol{A}^*$ 为 $\boldsymbol{A}$ 的伴随矩阵，若 $\boldsymbol{A}(\boldsymbol{A}-\boldsymbol{A}^*)=\boldsymbol{O}$ 且 $\boldsymbol{A}\neq\boldsymbol{A}^*$，则 $r(\boldsymbol{A})$ 取值为（　　）
\\
\noindent (A) 0或1. \qquad (B) 1或3. \qquad (C) 2或3. \qquad (D) 1或2.

\vfill
\newpage
\qn{2024.10}
设 $\boldsymbol{A},\boldsymbol{B}$ 为2阶矩阵，且 $\boldsymbol{AB}=\boldsymbol{BA}$，则“$\boldsymbol{A}$ 有两个不相等的特征值”是“$\boldsymbol{B}$ 可对角化”的（　　）
\\
\noindent (A) 充分必要条件. \qquad (B) 充分不必要条件.
\\
\noindent (C) 必要不充分条件. \qquad (D) 既不充分也不必要条件.

\vfill
\newpage
\qn{2024.11}
曲线 $y^2=x$ 在点 $(0,0)$ 处的曲率圆方程为________.

\vfill
\newpage
\qn{2024.12}
函数 $f(x,y)=2x^3-9x^2-6y^4+12x+24y$ 的极值点是________.

\vfill
\newpage
\qn{2024.13}
微分方程 $y'=\dfrac{1}{(x+y)^2}$ 满足条件 $y(1)=0$ 的解为________.

\vfill
\newpage
\qn{2024.14}
已知函数 $f(x)=x^2(\mathrm{e}^x+1)$，则 $f^{(5)}(1)=$________.

\vfill
\newpage
\qn{2024.15}
某物体以速度 $v(t)=t+k\sin\pi t$ 作直线运动. 若它从 $t=0$ 到 $t=3$ 的时间段内平均速度是 $\dfrac{5}{2}$，则 $k=$________.

\vfill
\newpage
\qn{2024.16}
设向量 $\boldsymbol{\alpha}_1=\begin{pmatrix}a\\1\\-1\\1\end{pmatrix}$，$\boldsymbol{\alpha}_2=\begin{pmatrix}1\\1\\b\\a\end{pmatrix}$，$\boldsymbol{\alpha}_3=\begin{pmatrix}1\\a\\-1\\1\end{pmatrix}$，若 $\boldsymbol{\alpha}_1,\boldsymbol{\alpha}_2,\boldsymbol{\alpha}_3$ 线性相关，且其中任意两个向量均线性无关，则 $ab=$________.

\vfill
\newpage
\qn{2024.17}
（本题满分10分）
设平面有界区域 $D$ 位于第一象限，由曲线 $xy=\dfrac{1}{3}$，$xy=3$ 与直线 $y=\dfrac{1}{3}x$，$y=3x$ 围成，计算 $\iint_D(1+x-y)\,\mathrm{d}x\mathrm{d}y$.

\vfill
\newpage
\qn{2024.18}
（本题满分12分）
设 $y(x)$ 为微分方程 $x^2y''+xy'-9y=0$ 的解，满足条件 $y\Big|_{x=1}=2$，$y'\Big|_{x=1}=6$.
\\
（1）利用变换 $x=\mathrm{e}^t$ 将上述方程化为常系数线性方程，并求 $y(x)$；
\\
（2）计算 $\int_1^2 y(x)\sqrt{4-x^2}\,\mathrm{d}x$.

\vfill
\newpage
\qn{2024.19}
（本题满分12分）
设 $t>0$，平面有界区域 $D$ 由曲线 $y=\sqrt{x}\,\mathrm{e}^{-x}$ 与直线 $x=t$，$x=2t$ 及 $x$ 轴围成，$D$ 绕 $x$ 轴旋转一周所成旋转体的体积为 $V(t)$，求 $V(t)$ 的最大值.

\vfill
\newpage
\qn{2024.20}
（本题满分12分）
已知函数 $f(u,v)$ 具有2阶连续偏导数，且函数 $g(x,y)=f(2x+y,3x-y)$ 满足 $\dfrac{\partial^2 g}{\partial x^2}+\dfrac{\partial^2 g}{\partial x\partial y}-6\dfrac{\partial^2 g}{\partial y^2}=1$.
\\
（1）求 $\dfrac{\partial^2 f}{\partial u\partial v}$；
\\
（2）若 $\dfrac{\partial f(u,0)}{\partial u}=u\mathrm{e}^{-u}$，$f(0,v)=\dfrac{1}{50}v^2-1$，求 $f(u,v)$ 的表达式.

\vfill
\newpage
\qn{2024.21}
（本题满分12分）
设函数 $f(x)$ 具有2阶导数，且 $f'(0)=f'(1)$，$|f''(x)|\leqslant1$. 证明：
\\
（1）当 $x\in(0,1)$ 时，$|f(x)-f(0)(1-x)-f(1)x|\leqslant\dfrac{x(1-x)}{2}$；
\\
（2）$\left|\int_0^1f(x)\,\mathrm{d}x-\dfrac{f(0)+f(1)}{2}\right|\leqslant\dfrac{1}{12}$.

\vfill
\newpage
\qn{2024.22}
（本题满分12分）
设矩阵 $\boldsymbol{A}=\begin{pmatrix}0&1&a\\1&0&1\end{pmatrix}$，$\boldsymbol{B}=\begin{pmatrix}1&1\\1&1\\b&2\end{pmatrix}$，二次型 $f(x_1,x_2,x_3)=\boldsymbol{x}^{\mathrm{T}}\boldsymbol{B}\boldsymbol{A}\boldsymbol{x}$. 已知方程组 $\boldsymbol{Ax}=\boldsymbol{0}$ 的解均是 $\boldsymbol{B}^{\mathrm{T}}\boldsymbol{x}=\boldsymbol{0}$ 的解，但这两个方程组不同解.
\\
（1）求 $a,b$ 的值；
\\
（2）求正交变换 $\boldsymbol{x}=\boldsymbol{Qy}$ 将 $f(x_1,x_2,x_3)$ 化为标准形.

\vfill
\newpage
\end{document}